\documentclass[11pt]{article}
\usepackage[margin=1in]{geometry}
\usepackage{amsmath,amssymb,amsthm,hyperref}
\newtheorem{lemma}{Lemma}\newtheorem{proposition}[lemma]{Proposition}
\title{The primitive physical moment pair and periodic local expansions}
\author{Internal research note}\date{18 November 2026}
\begin{document}\maketitle
For any prime $p$, put $q=p-1$ and, for $n\ge1$,
\[
 H_n(Y)=p\sum_{k=0}^n\binom nk(n)_{qk}(Y/p)^k,\quad
 U_p(k)=\frac{(pk)!}{p^kk!}
       =\prod_{\substack{1\le j\le pk\\p\nmid j}}j,\quad U_p(0)=1.
\]
The physical coefficient identity is
\[
 [Y^k]H_n=U_p(k)
 \left[p\binom{n+qk-1}{pk}+\binom{n+qk-1}{pk-1}\right],
 \quad k\ge1.
 \tag{1}
\]
We use the formal convention $H_0=p$ only in generating series;
no Gamma distribution at shape zero is asserted.
The generating series below are formal in $z$, with polynomial
coefficients in $Y$.

\begin{lemma}[The primitive pair, valid at every prime]\label{lem:pair}
For $n\ge1$,
\[
 W_n(Y)=\frac{nH_n-YH'_n}{pn}
       =\sum_{k=0}^{n-1}U_p(k)
                  \binom{n+qk-1}{pk}Y^k\in\mathbb Z[Y],
 \quad W_n(0)=1,\quad [Y^{n-1}]W_n=U_p(n-1),
 \tag{2}
\]
and
\[
                   H_n=W_{n+1}+qW_n.                    \tag{3}
\]
Put $Q=Yz/(1-z)^p$ and $M_p(Q)=\sum_{k\ge0}U_p(k)Q^k$.
Then
\[
 G_p(Y,z):=\sum_{n\ge1}W_n(Y)z^{n-1}
       =\frac{M_p(Q)}{1-z},\qquad
 \sum_{n\ge0}H_n(Y)z^n=q+(1+qz)G_p(Y,z).
 \tag{4}
\]
\end{lemma}
\begin{proof}
For $k<n$, direct cancellation in the factorial expression for
$[Y^k]H_n$ gives
\[
 \frac{n-k}{pn}[Y^k]H_n
 =U_p(k)\binom{n+qk-1}{pk}.
\]
The $k=n$ term vanishes, and $k=0$ gives one.
Pascal's identity gives coefficientwise
\[
 U_p(k)\left[\binom{n+qk}{pk}
          +q\binom{n+qk-1}{pk}\right]=[Y^k]H_n,
\]
including the endpoints. This proves (2)--(3).
Writing $n=k+j$, $j\ge1$, the elementary binomial series gives
\[
 \sum_{j\ge1}\binom{pk+j-1}{pk}z^{j-1}
                      =(1-z)^{-pk-1}.
\]
Sum over $k$. The second identity in (4) follows from (3)
and $W_1=1$, including the constant convention $H_0=p$.
The coefficient at degree $n-1$ is $U_p(n-1)$ because the
binomial coefficient there equals one.
\end{proof}


\begin{corollary}[The actual consecutive moment pair]\label{cor:critical}
For $n\ge2$ and $Y=pt$ with $t=B/A$, the ordinary factorial
moments of $X(A+BX^{p-1})$ at $n-1,n$, divided respectively by
$A^{n-1}(n-1)!$ and $A^n n!$, are $W_n(Y), W_{n+1}(Y)$.
This formula uses $A\ne0$. Their common roots are the
multiple roots of $H_n$ and have valuation zero over every
extension of $\mathbb Q_p$.
\end{corollary}
\begin{proof}
Write $T(t)=H_n(pt)/p$ and $\theta=tD_t$.
In the lower moment each binomial coefficient is multiplied
by $(n-k)/n$, so the normalized polynomial is $T-\theta T/n$.
In the upper moment the rising factorial changes from
$(n)_{qk}$ to $(n+1)_{qk}=(1+qk/n)(n)_{qk}$.
The normalized polynomial is $T+q\theta T/n=H_n-qW_n=W_{n+1}$.
Equation (2) gives the equivalence with multiplicity because
$H_n(0)=p$. A root of $W_n$ of positive valuation would have a
unique smallest term, its constant term. For a root of negative
valuation the leading term would be uniquely smallest.
Both valuations are impossible.
\end{proof}

\begin{lemma}[The unit-factorial period multiplier]\label{lem:antiperiod}
For every prime $p$ and $K\ge1$, put $L=p^{K-1}$ and
\[
 \epsilon_{p,K}=\begin{cases}
 1,&p=2,\ K\ge3,\\
 -1,&\text{otherwise}.
 \end{cases}
\]
Then $U_p(k+L)\equiv\epsilon_{p,K}U_p(k)\pmod{p^K}$
for all $k\ge0$. With $A_{p,K}(Q)=\sum_{j=0}^{L-1}U_p(j)Q^j$,
\[
 M_p(Q)\equiv\frac{A_{p,K}(Q)}{1-\epsilon_{p,K}Q^L}\pmod{p^K}.
 \tag{5}
\]
\end{lemma}
\begin{proof}
The ratio is the product of every unit residue in an interval
of length $p^K$. Pair each residue with its inverse. For odd
$p$ only $1$ and $-1$ are self-inverse, since the factors of
$(x-1)(x+1)$ cannot both be divisible by $p$.
For $p=2$, the classes are $1$ at $K=1$ and $\{\pm1\}$ at $K=2$.
For $K\ge3$, one of $x-1,x+1$ for odd $x$ has valuation exactly
one, so the other must be divisible by $2^{K-1}$.
The four solutions are $\pm1$ and $\pm1+2^{K-1}$, whose product
is one modulo $2^K$. This proves the multiplier in every case.
Summing by the remainder of $k$ modulo $L$ gives (5).
\end{proof}

\begin{proposition}[A compatible rational recurrence at every precision]
\label{prop:rational}
Under the hypotheses of Lemma~\ref{lem:antiperiod}, put
\[
 \begin{split}
 D_{p,K}(Y,z)&=(1-z)^{pL}-\epsilon_{p,K}(Yz)^L,\\
 R_{p,K}(Y,z)&=
 \sum_{j=0}^{L-1}U_p(j)(Yz)^j(1-z)^{p(L-j)-1}.
 \end{split}
\]
Coefficientwise modulo $p^K$,
\[
 G_p=\frac{R_{p,K}}{D_{p,K}},\qquad
 \sum_{n\ge0}H_nz^n=q+(1+qz)\frac{R_{p,K}}{D_{p,K}}.
 \tag{7}
\]
Here $\deg_zD_{p,K}=p^K$, its constant coefficient is one,
its leading coefficient is $(-1)^{pL}$, and
$\deg_zR_{p,K}\le p^K-1$.
Both actual sequences therefore satisfy a reversible linear
recurrence of order $p^K$ from index $n\ge p^K+1$.
For a fixed specialization $Y$ in any finite unramified
coefficient ring $\mathcal O_f/p^K$, the sequence $W_n(Y)$
is purely periodic: its length-$p^K$ state is advanced by an
invertible companion matrix in a finite ring.
This is finite precision at a fixed coefficient ring, not a
uniform finite state space for all precisions and residue degrees.
\end{proposition}
\begin{proof}
Substitute (6) into (4) and clear $(1-z)^{pL}$.
Every exponent in $R_{p,K}$ is nonnegative.
The highest possible degree is $pL-1$.
The leading coefficient of $(1-z)^{pL}$ is a unit. The additional term has smaller degree.
These are precisely the asserted denominator and recurrence
properties. Compatibility between precisions follows because
each expression is congruent to the same exact series (4).
\end{proof}

\begin{lemma}[Uniform cubic precision for $p\ge5$]\label{lem:cubic}
Let $p\ge5$ and $w=(p-1)!$. For every $k\ge0$,
\[
                        U_p(k)\equiv w^k\pmod{p^3}.       \tag{8}
\]
In particular the entire actual moment pair satisfies
\[
 G_p(Y,z)\equiv\frac{(1-z)^{p-1}}{(1-z)^p-wYz}\pmod{p^3},
 \quad
 H_n(Y)\equiv K_{n,p}(wY)\pmod{p^3},                     \tag{9}
\]
where $K_{n,p}(Z)$ is the power sum of the $p$ roots of
$(X-1)^p-ZX^{p-1}$. These are all-index and
coefficientwise congruences, with no restriction on $k$ or $n$
relative to $p$.
The cubic assertion fails at $p=3$:
$U_3(2)=40\not\equiv 2^2\pmod {27}$.
\end{lemma}
\begin{proof}
Set $S_1=\sum_{i=1}^{p-1}i^{-1}$ and
$S_2=\sum_{i=1}^{p-1}i^{-2}$ in $\mathbb Z_p$.
Multiplication by a residue $b$ with $b^2\ne1$ shows
$S_2=0\bmod p$. Pairing $i$ and $p-i$ gives
\[
 S_1=\frac12\sum_i\left(\frac1i+\frac1{p-i}\right)
           \equiv-\frac p2S_2\equiv0\pmod {p^2}.
\]
For any integer $\ell\ge0$,
\[
 \prod_{i=1}^{p-1}(p\ell+i)
 \equiv w\left(1+p\ell S_1+
          \frac{p^2\ell^2}{2}(S_1^2-S_2)\right)
 \equiv w\pmod {p^3}.
\]
Multiplying the finitely many blocks proves (8), for every $k$.
Thus $M_p(Q)\equiv(1-wQ)^{-1}\bmod p^3$, proving the
first formula in (9) using (4).
The second follows from the exact rational power-sum identity
\[
 \sum_{n\ge0}K_{n,p}(wY)z^n
  =p-\frac{zD'(z)}{D(z)}
  =q+(1+qz)\frac{(1-z)^{p-1}}{D(z)},\qquad
 D(z)=(1-z)^p-wYz.
\]
The numerical $p=3$ control proves the stated scope limitation.
\end{proof}

\begin{theorem}[Explicit periods for truncated local expansions]\label{thm:jet-period}
Fix a prime $p$, a precision $K\ge1$, and a Taylor order $e\ge0$.
Let $y$ be a unit in the integers $\mathcal O_E$ of a finite
unramified extension $E/\mathbb Q_p$. Choose $d$ such that
$\mathbb F_{p^d}$ contains the residue of $y$ and all roots of
\[
 f_{\bar y}(X)=X^p+\bar y X^{p-1}-1 .
\]
Let $m=p^d-1$, and let $t_e$ be the least nonnegative integer
with $p^{t_e}>e$. Modulo $(p^K,Z^{e+1})$ the two sequences
$H_n(y+Z)$ and $W_n(y+Z)$, $n\ge1$, have common period
\[
 P=m p^{K-1+\max(K-1,t_e)} .
 \tag{10}
\]
If $p\ge5$ and $K\le3$, the smaller bound $m p^{K-1+t_e}$
also gives a period.
\end{theorem}
\begin{proof}
Write $L=p^{K-1}$.
The reciprocal of the denominator of (7) is monic of degree
$pL$, with unit constant term. Modulo $p$ it equals
$f_{\bar y+Z}(X)^L$. In characteristic $p$,
$f_{\bar y}$ has simple nonzero roots because its derivative is
$-\bar y X^{p-2}$ and its constant term is $-1$.
Thus $f_{\bar y}$ divides $X^m-1$ in the residue field
of $E$. In the reciprocal quotient modulo $p$ the element
$X^m-1$ lies in the ideal $(f_{\bar y+Z},Z)$, since
$f_{\bar y}=f_{\bar y+Z}-ZX^{p-1}$.
With $h=\max(K-1,t_e)$, Frobenius to exponent $p^h$
annihilates that element. Indeed it kills both summands, as
$p^h\ge L$ and $p^h>e$. Hence $X^{mp^h}=1+pB$ in the quotient algebra.
Taking the $p^{K-1}$st power gives $X^P=1$ modulo $p^K$.
The reciprocal polynomial annihilates the companion shift of
the recurrence for $W_n$, from the first full length-$pL$ state.
Therefore the shift to the $P$th power is the identity, proving
periodicity from $n=1$. Equation (3) gives the same period for $H_n$.

For $p\ge5$ and $K\le3$, use the denominator of (9) instead.
Its reciprocal reduces to $f_{\bar y+Z}$, since $w\equiv-1\bmod p$.
Thus the same argument uses $L=1$ at the residue step and
$h=t_e$. The numerator in (9) has degree below the denominator
degree, so this recurrence also starts with the first full state.
\end{proof}

\begin{corollary}[A local exclusion in an arithmetic progression]\label{cor:145}
For every positive integer $n\equiv145\pmod{16758}$, the polynomial
$H_{n,7}(Y)$ has no multiple root in the full residue disk
$v_7(Y-4)>0$. Its roots there have distances of valuations $1/5$
and $1/9$ from $4$, in groups of five and nine. Both groups give
Eisenstein factors over $\mathbb Q_7$ after translation by $4$.
\end{corollary}
\begin{proof}
At $p=7$ and $\bar y=4$, an inverse root $\tau$ of the reciprocal
equation satisfies $\tau^7=1+4\tau$.
Three iterations give $\tau^{343}=\tau$ since
$4^3=1$ and $1+4+4^2=0$ in $\mathbb F_7$.
In Theorem~\ref{thm:jet-period} take $d=3$ and $m=342$.
For the three pairs $(K,e)=(3,0),(2,5),(1,14)$, the stronger bound
gives the same period $342\cdot7^2=16758$.

Write $H_{145,7}(4+Z)=\sum a_j Z^j$.
The coefficient formula (1) and the binomial expansion give
\[
 (a_j\bmod343)_{0\le j\le14}=
 (98,147,245,98,196,336,175,252,252,238,245,315,49,294,250).
\]
In particular, the displayed precision levels certify $v_7(a_0)=2$,
$v_7(a_j)\ge2$ for $1\le j\le4$, $v_7(a_5)=1$,
$v_7(a_j)\ge1$ for $6\le j\le13$, and $v_7(a_{14})=0$.
These inequalities and equalities therefore hold for the coefficients
of $H_{n,7}(4+Z)$ at every index in the progression.
All coefficients above degree $14$ are integral by (1).

For a positive weight $s$, the minimum of $v_7(a_j)+js$ is therefore
$\min(2,1+5s,14s)$.
This weighted valuation is multiplicative on polynomials, since
the weighted initial terms multiply in the graded domain of the
valued ring. Factoring over an algebraic closure shows that each
root $\alpha$ of the translated polynomial contributes
$\min(s,v_7(\alpha))$ to the weighted valuation.
Its slope drops by five at $s=1/5$ and by nine at $s=1/9$.
Thus the two groups are exactly the roots of positive valuation.
Each group is Galois invariant and its monic factor has constant
valuation one, with all other nonleading coefficients of positive
integer valuation. Each factor is Eisenstein over $\mathbb Q_7$,
hence separable, and the two factors have disjoint roots.
\end{proof}

The progression concerns the residue disk at $4$.
Excluding the remaining residue disks is required for an
all-index squarefreeness proof.
\end{document}
